\subsection{Theoretical context}

The formal properties of the doubly weighted statistic $\hat\theta_{\mathrm{both}}$ follow from existing theory applied to this estimand. We state the regularity under which the main-text sketches hold; full derivations are standard applications of cited results.

\textbf{Equivalence to weighted AUC.} The Mann-Whitney functional $\Psi$ is equivalent to the area under the ROC curve~\cite{Hanley_1982}. The weighted score distributions in Section~\ref{sec:method} are therefore weighted AUC statistics. This equivalence is exact for the empirical distributions and carries to the population quantities $\Theta$ and $\Theta_\lambda$ defined in Section~\ref{sec:method}.

\textbf{U-statistic structure and Hadamard conditions.} The Mann-Whitney statistic is a degree-2 $U$-statistic~\cite{Lee_1990}. With additive importance weights applied per observation, the weighted version retains $U$-statistic structure. O'Neil and Redner~\cite{ONeil_1993} established asymptotic normality for weighted $U$-statistics of degree~2 under mild moment conditions on the weights (bounded second moments of $w_{\mathrm{source}}$ and $w_{\mathrm{target}}$ and non-degenerate kernel variance). Lee~\cite{Lee_1990} provides variance formulas via the H\'ajek projection and Hoeffding decomposition. For Proposition~\ref{prop:consistency}, Hadamard differentiability of $\Psi$ requires absolutely continuous score CDFs with bounded densities (so the ROC curve is continuously differentiable), and Hadamard differentiability of the weighting map $Q \mapsto Q_\lambda$ requires $\hat{r} \to_p r$ in $L^2(P_S)$ with $\sup_x \hat{r}(x)$ bounded in probability and $\lambda \in [0,1]$ fixed; the composition then inherits differentiability via the functional delta method and the continuous mapping theorem.

\textbf{Weighted MWW generalization.} The weighted Mann-Whitney-Wilcoxon statistic~\cite{Xie_2002} includes our test statistic as a special case with RIW density-ratio weights. Permutation validity (Proposition~\ref{prop:permutation}) additionally requires that weights be a function of features $x$ only and held fixed across permutations; weight estimation error then affects the interpretation of $\overlap$ but not conditional Type~I control given the weight configuration~\cite{Lehmann_2005}.

\subsection{Bias of the $\lambda$-Interpolated Estimand}

For $\lambda > 0$, $\Theta_\lambda \neq \Theta$ because the $\lambda$-interpolated weights only partially correct the density ratio. Under $L^2$ convergence of $\hat{r}$ to $r$ and bounded total variation distance $d_{\mathrm{TV}}(P_S,P_T)$, the bias satisfies $\|\Theta_\lambda - \Theta\| = O(1-\lambda) \cdot d_{\mathrm{TV}}(P_S,P_T)$: at $\lambda=0$ the weights implement exact density-ratio correction ($\Theta_0=\Theta$), at $\lambda=1$ they are uniform ($\Theta_1$ is the unweighted population quantity), and intermediate $\lambda$ interpolates linearly in the weight space. Smaller $\lambda$ therefore reduces bias toward the common-support estimand but increases variance through weight concentration (lower $\min(\mathrm{ESS}_{\mathrm{src}},\mathrm{ESS}_{\mathrm{tgt}})$). This is a statement about the population interpolation, not about finite-sample variance of $\hat{r}$; practitioners should consult the $\lambda$ sensitivity curve and ESS diagnostic jointly.
